\chapter{CONCLUSION AND FUTURE WORK}

\section{Quantum Machine Learning Enhancements}

Future improvements to the system can enhance quantum feature representation and classification performance.

\subsection{Expand Qubit Count}

Increasing the number of qubits from 4 to 6-8 would expand the feature space to 64-256 dimensions.

Maintaining PCA variance between 80-90\% may improve kernel expressiveness and reduce support vector ratios.

\subsection{Advanced Quantum Feature Maps}

Future implementations may replace the existing feature map with advanced circuits such as:

\begin{itemize}
\item Pauli Feature Maps
\item IQP circuits
\end{itemize}

Trainable quantum kernels can also be introduced to maximize class separability.

\subsection{Hybrid Variational Quantum Circuits}

Variational Quantum Eigensolver (VQE) or Quantum Neural Networks may be integrated before the SVM classifier to learn latent quantum representations.

\section{Classical Machine Learning Enhancements}

\subsection{Improved Label Design}

Future work may incorporate physics-based thresholds such as State of Health (SOH) limits instead of statistical quantile labels.

\subsection{Enhanced Feature Engineering}

Additional electrochemical features such as Electrochemical Impedance Spectroscopy (EIS) parameters may be included.

Wavelet transforms can also be used to extract multi-frequency degradation signatures.

\subsection{Classical Baseline Comparison}

Benchmarking against classical models such as:

\begin{itemize}
\item Random Forest
\item XGBoost
\item LSTM
\end{itemize}

can help evaluate the practical advantages of quantum models.

\section{Data and Scalability}

Future experiments may include larger datasets from NASA and CALCE battery repositories to improve generalization.

Testing on real quantum hardware using IBM Quantum devices can also evaluate practical quantum advantage.

\section{System Integration and Deployment}

The risk prediction module can be integrated into Battery Management Systems (BMS) as a REST API microservice.

Real-time dashboards can visualize battery risk levels and trigger alerts when critical thresholds are reached.

Explainability methods such as SHAP can be used to interpret feature contributions in the model predictions.

\section{Multi-Chemistry Extension}

Future systems may extend the methodology to other battery chemistries such as lithium-sulfur or solid-state batteries, enabling broader battery safety monitoring applications.